
\documentclass{article}
\usepackage{colortbl}
\usepackage{makecell}
\usepackage{multirow}
\usepackage{supertabular}

\begin{document}

\newcounter{utterance}

\centering \large Interaction Transcript for game `cladder', experiment `full\_v1.5\_default', episode 6302 with Qwen3.6{-}27B.
\vspace{24pt}

{ \footnotesize  \setcounter{utterance}{1}
\setlength{\tabcolsep}{0pt}
\begin{supertabular}{c@{$\;$}|p{.15\linewidth}@{}p{.15\linewidth}p{.15\linewidth}p{.15\linewidth}p{.15\linewidth}p{.15\linewidth}}
    \# & \multicolumn{2}{c}{Player} && \multicolumn{2}{c}{Game Master} \\
    \hline

    \theutterance \stepcounter{utterance}  
    & & & \multicolumn{4}{p{0.6\linewidth}}{
        \cellcolor[rgb]{0.9,0.9,0.9}{
            \makecell[{{p{\linewidth}}}]{
                \texttt{\tiny{[P1$\langle$GM]}}
                \texttt{You are an expert in causal inference. The following question is not a typical commonsense query, but rather a meticulously designed question created by a professor specializing in causal inference, intended to assess the students' mastery of the course content.} \\
\texttt{Imagine a self{-}contained, hypothetical world with only the following conditions, and without any unmentioned factors or causal relationships: Zuph has a direct effect on rixq and xevu. Rixq has a direct effect on xevu. For those who are not zuph and are not rixq, the probability of xevu is 15\%. For those who are not zuph and are rixq, the probability of xevu is 50\%. For those who are zuph and are not rixq, the probability of xevu is 32\%. For those who are zuph and are rixq, the probability of xevu is 85\%. The overall probability of zuph is 84\%. Will rixq increase the chance of xevu?} \\
\texttt{Start your answer with 'Yes' or 'No', followed by additional reasoning or evidence to support your explanation.} \\
            }
        }
    }
     \\ \\

    \theutterance \stepcounter{utterance}  
    & \multicolumn{4}{p{0.6\linewidth}}{
        \cellcolor[rgb]{0.561,0.737,0.561}{
            \makecell[{{p{\linewidth}}}]{
                \texttt{\tiny{[P1$\rangle$GM]}}
                \texttt{Yes} \\
\\ 
\texttt{To determine whether `rixq` increases the chance of `xevu`, we must calculate the Average Causal Effect (ACE) of `rixq` on `xevu`. In causal inference, this is defined as the difference between the expected outcome if everyone were treated (`rixq = 1`) and the expected outcome if no one were treated (`rixq = 0`).} \\
\\ 
\texttt{The problem provides the conditional probabilities of `xevu` given specific combinations of `zuph` and `rixq`. Let \$P(X|Z, R)\$ denote the probability of `xevu` given `zuph` (\$Z\$) and `rixq` (\$R\$). The given values are:} \\
\texttt{1.  \$P(X|Z=0, R=0) = 0.15\$} \\
\texttt{2.  \$P(X|Z=0, R=1) = 0.50\$} \\
\texttt{3.  \$P(X|Z=1, R=0) = 0.32\$} \\
\texttt{4.  \$P(X|Z=1, R=1) = 0.85\$} \\
\\ 
\texttt{We are also given the marginal probability of `zuph`:} \\
\texttt{\$P(Z=1) = 0.84\$} \\
\texttt{Therefore, \$P(Z=0) = 1 {-} 0.84 = 0.16\$.} \\
\\ 
\texttt{The causal effect of} \\
            }
        }
    }
    & & \\ \\

    \theutterance \stepcounter{utterance}  
    & & & \multicolumn{2}{p{0.3\linewidth}}{
        \cellcolor[rgb]{0.9,0.9,0.9}{
            \makecell[{{p{\linewidth}}}]{
                \texttt{\tiny{[GM$|$GM]}}
                \texttt{Yes} \\
            }
        }
    }
    & & \\ \\

    \theutterance \stepcounter{utterance}  
    & & & \multicolumn{2}{p{0.3\linewidth}}{
        \cellcolor[rgb]{0.9,0.9,0.9}{
            \makecell[{{p{\linewidth}}}]{
                \texttt{\tiny{[GM$|$GM]}}
                \texttt{yes} \\
            }
        }
    }
    & & \\ \\

    \theutterance \stepcounter{utterance}  
    & & & \multicolumn{2}{p{0.3\linewidth}}{
        \cellcolor[rgb]{0.9,0.9,0.9}{
            \makecell[{{p{\linewidth}}}]{
                \texttt{\tiny{[GM$|$GM]}}
                \texttt{game\_result = WIN} \\
            }
        }
    }
    & & \\ \\

\end{supertabular}
}

\end{document}
